% ============================================================
% Come si usano, argomento 2: calcolo combinatorio
% ============================================================
\argomento{Calcolo combinatorio}{Formule a p.~\pageref*{f:fondamentale}}

\begin{scelta}{Come scegliere la formula: due domande}
Prima di contare ci si chiede: \textbf{l'ordine conta?} (si contano sequenze o gruppi?) e \textbf{gli elementi si
possono ripetere?}
\begin{center}
\renewcommand{\arraystretch}{1.7}
\begin{tabular}{@{}r|c|c@{}}
  & \textbf{ripetizioni sì} & \textbf{ripetizioni no} \\ \hline
  \textbf{ordine sì} & $n^k$\enspace\vv{disprip}
    & $\dfrac{n!}{(n-k)!}$\enspace\vv{dispsemplici}\qquad con $k = n$: \ $n!$\enspace\vv{permutazioni} \\ \hline
  \textbf{ordine no} & $\dbinom{n+k-1}{k}$\enspace\vv{combrip} & $\dbinom{n}{k}$\enspace\vv{combinazioni} \\
\end{tabular}
\end{center}
\begin{itemize}[leftmargin=1.2em, itemsep=2pt, topsep=0pt]
  \item \textbf{L'ordine conta} se le posizioni sono diverse fra loro: podi, PIN, parole, sequenze di lanci,
        estrazioni in cui si registra chi esce per primo. \textbf{Non conta} se interessa solo \emph{chi} c'è: mani
        di carte, gruppi, comitati, schedine.
  \item \textbf{Si ripete} se un elemento può ricomparire: lanci di un dado, cifre di un PIN, estrazioni
        \emph{con} reinserimento. \textbf{Non si ripete} nelle estrazioni \emph{senza} reinserimento.
  \item Se la scelta si fa in passi \textbf{diversi} (le lettere e le cifre di una targa; prima gli assi, poi le
        altre carte) si conta ogni passo e si \textbf{moltiplica} (\vv{fondamentale}). Se il testo distingue casi
        che non si sovrappongono («esattamente tre oppure esattamente quattro assi») si conta ogni caso e si
        \textbf{somma}.
\end{itemize}
\end{scelta}

\scheda{fondamentale}{
  \usoquando{La scelta si fa in passi successivi e a ogni passo il \emph{numero} di possibilità è sempre lo stesso,
    qualunque siano le scelte fatte prima (le possibilità possono cambiare, il loro numero no).}
  \usocome{Si elencano i passi, si scrive sotto ciascuno quante scelte offre e si moltiplica. Un diagramma ad
    albero aiuta: le foglie sono i risultati.}
  \usoattenzione{Se il numero di scelte di un passo dipende da quelle precedenti (per esempio «il secondo numero è
    maggiore del primo»), il prodotto non si può fare direttamente: si divide in casi e si sommano.}
  \usoesempio{Targhe \texttt{AB 123 CD} con 22 lettere ammesse:
    $22 \cdot 22 \cdot 10 \cdot 10 \cdot 10 \cdot 22 \cdot 22 = 22^4 \cdot 10^3 = 234\,256\,000$.}
}

\scheda{disprip}{
  \usoquando{Sequenze di $k$ simboli scelti fra $n$, in cui conta la posizione e un simbolo può ripetersi: PIN,
    parole, stringhe binarie, $k$ lanci di un dado, estrazioni con reinserimento.}
  \usocome{$n$ = quanti simboli sono disponibili a ogni posizione; $k$ = quante posizioni.}
  \usoattenzione{È $n^k$, non $k^n$: la base è il numero di simboli, l'esponente la lunghezza.}
  \usoesempio{PIN di 4 cifre: $10^4 = 10\,000$. Un byte: $2^8 = 256$ valori. Tre lanci di un dado:
    $6^3 = 216$ risultati.}
}

\scheda{dispsemplici}{
  \usoquando{$k$ oggetti distinti da mettere in $k$ posizioni diverse: podi, cariche (presidente, vice,
    segretario), estrazioni senza reinserimento in cui conta l'ordine di uscita.}
  \usocome{Si moltiplicano $k$ fattori che scendono di uno alla volta partendo da $n$. Con i fattoriali:
    $n!/(n-k)!$.}
  \usoesempio{Otto atleti in finale, podi possibili: $8 \cdot 7 \cdot 6 = \frac{8!}{5!} = 336$. Quaterne ordinate
    da un mazzo di 52 carte: $52 \cdot 51 \cdot 50 \cdot 49 = 6\,497\,400$.}
}

\scheda{permutazioni}{
  \usoquando{Si mettono in fila \emph{tutti} gli $n$ oggetti: anagrammi, ordini di arrivo, modi di mescolare un
    mazzo.}
  \usovarianti{Se alcuni oggetti sono uguali, scambiarli fra loro non produce sequenze nuove: si divide per i
    fattoriali delle ripetizioni. È lo stesso ragionamento delle stringhe con $k$ uni (\vv{stringhe}).}
  \usoesempio{Anagrammi di ROMA: $4! = 24$. Anagrammi di ANNA (due A e due N): $\frac{4!}{2!\,2!} = 6$.}
}

\scheda{fattoriale}{
  \usoquando{In ogni formula di conteggio; il problema è solo fare i conti senza numeri enormi.}
  \usocome{Nei rapporti si cancellano i fattori comuni prima di moltiplicare: $\frac{n!}{m!}$ è il prodotto dei
    numeri da $m + 1$ a $n$.}
  \usoattenzione{$0! = 1$ (c'è un solo modo di mettere in fila zero oggetti). Serve, per esempio, in
    $\binom{n}{0} = \frac{n!}{0!\,n!} = 1$.}
  \usoesempio{$\frac{7!}{4!} = 7 \cdot 6 \cdot 5 = 210$; \ $5! = 5 \cdot 4! = 120$.}
}

\scheda{combinazioni}{
  \usoquando{Si scelgono $k$ elementi distinti fra $n$ e interessa solo \emph{quali}, non in che ordine: mani di
    carte, gruppi, comitati, numeri giocati, posizioni dei successi in una sequenza.}
  \usocome{Si contano le disposizioni e si divide per $k!$, perché ogni gruppo è stato contato in tutti i suoi
    $k!$ ordini: $k$ fattori decrescenti sopra, $k!$ sotto.}
  \usovarianti{Se il gruppo deve contenere un numero preciso di elementi di tipi diversi («esattamente due assi»),
    si sceglie separatamente da ogni tipo e si moltiplica.}
  \usoattenzione{«Esattamente» e «almeno» sono domande diverse: per «almeno tre assi» si sommano i casi «tre» e
    «quattro», che non si sovrappongono.}
  \usoesempio{Giocate del SuperEnalotto: $\binom{90}{6} = 622\,614\,630$. Cinquine da 52 carte con esattamente due
    assi: $\binom{4}{2}\binom{48}{3} = 6 \cdot 17\,296 = 103\,776$.}
}

\scheda{binomiale}{
  \usoquando{Per accorciare i conti e per controllarli.}
  \usocome{Con la simmetria si usa il $k$ più piccolo: $\binom{90}{84} = \binom{90}{6}$. Con la regola di Stifel
    si costruisce il triangolo di Tartaglia: ogni numero è la somma dei due che ha sopra.}
  \usoesempio{$\binom{8}{5} = \binom{8}{3} = \frac{8 \cdot 7 \cdot 6}{3!} = 56$. \ Stifel:
    $\binom{5}{2} = \binom{4}{1} + \binom{4}{2} = 4 + 6 = 10$.}
}

\scheda{combrip}{
  \usoquando{Si scelgono $k$ oggetti fra $n$ tipi, i tipi si possono ripetere e l'ordine non conta.}
  \usocome{Si disegnano $k$ pallini (gli oggetti) e $n - 1$ stanghette che separano i tipi: ogni scelta è una
    stringa di $n + k - 1$ simboli, in cui si decide dove stanno i $k$ pallini.}
  \usoesempio{Tre palline di gelato fra cinque gusti: $\binom{5 + 3 - 1}{3} = \binom{7}{3} = 35$.}
}

\scheda{stringhe}{
  \usoquando{Si contano le sequenze di $n$ prove (bit, lanci) con esattamente $k$ successi. È il primo
    ingrediente della binomiale (\vv{distbinomiale}).}
  \usocome{Una sequenza è individuata dalle posizioni dei suoi $k$ successi: si sceglie un sottoinsieme di $k$
    posizioni fra $n$.}
  \usoesempio{Dieci lanci di una moneta con esattamente tre teste: $\binom{10}{3} = 120$ sequenze.}
}

\scheda{newton}{
  \usoquando{Per sviluppare una potenza di un binomio e per sommare coefficienti binomiali. È il motivo per cui la
    pmf binomiale somma a 1.}
  \usoesempio{$\sum_{k=0}^{N} \binom{N}{k}\, p^k (1-p)^{N-k} = (p + 1 - p)^N = 1$.}
}

\scheda{sottoinsiemi}{
  \usoquando{Per contare tutti i sottoinsiemi di un insieme finito, di qualunque dimensione: per esempio tutti gli
    eventi di uno spazio campione finito.}
  \usocome{Ogni sottoinsieme è una stringa binaria (1 = l'elemento c'è): le stringhe lunghe $n$ sono $2^n$.}
  \usoesempio{Gli eventi del lancio di un dado sono $2^6 = 64$, compresi $\emptyset$ e $\Omega$.}
}
